\newcommand{\mainFindings}{Convex and Homotopy have the lowest mean patch RMSE (both about 0.845 mm/h). ILDW-init Roy has mean RMSE 1.216 mm/h versus 1.239 for ILDW, while GMZ has 1.263 mm/h. Roy fits real-link attenuation most closely, but this does not translate into the lowest rainfall-map error.}
\newcommand{\rainFindings}{The mean patch RMSE reduction from ILDW to ILDW-init Roy is 1.86\%. GMZ's mean RMSE is 1.92\% higher than ILDW's. Convex and Homotopy are very close at the reported precision: their mean RMSEs are 0.845199 and 0.845239 mm/h, respectively. The difference between mean and median RMSE reflects a tail of difficult patches; the plot retains all of them.}
\newcommand{\attenFindings}{Roy has the smallest mean link-normalized error and the smallest length-weighted error in this seven-method comparison. The map-error ranking is different: Convex and Homotopy lead the rainfall RMSE table. Link measurements constrain path integrals, so close attenuation agreement still allows different spatial rainfall patterns.}
\newcommand{\referenceSD}{3.130}
\newcommand{\royFindings}{The diagnostics report 52/100 successful outer stops; the recorded reasons are 48 max outer, 52 local fixed point. Median outer iterations are 297. Across all patches, 24,794 of 24,794 inner calls report success, with 3,593,795 ADMM iterations and 6,906,467 CG iterations. The median final objective ratio is 0.0909; the median final full-step/tolerance ratio is 0.999. Convergence of the inner solves does not establish convergence of every nonlinear outer run.}
\newcommand{\gmzFindings}{There are 19,741--50,978 gauges per patch (median 29,892). The median logged virtual-gauge attenuation RMSE is 2.10e-15 dB, while the median final-raster attenuation RMSE is 0.0440 dB. The number of gauges without other-link neighbors ranges from 0 to 0 per patch.}
